\section{Introduction}

The hard problem of consciousness (Chalmers, 1995) names the gap between physical process and felt experience. Most responses either deny that experience is real or treat it as an extra property layered onto an already physical world. This paper takes the opposite route. It treats experience as the only substance there is, and recovers the physical as a large-scale pattern within it.

This is the second presentation of Vibe Theory. The first (Pollard, 2025) set out the framework as philosophy. This paper gives it a discrete mathematical structure and commits to a specific minimal model. The aim is not yet empirical prediction but a coherent, simulable model whose pieces are all forced, derived, or constructed from a single axiom and two constraints.

The framework is panpsychist in spirit, building on the recent monism of Strawson (2006) and Goff (2017), and on the process philosophy of Whitehead (1929) and Russell (1927). It treats experience as ontologically primitive, as Integrated Information Theory treats consciousness as fundamental (Tononi, 2008), but with a different mathematical form. Where it departs from earlier panpsychism is in the strictness of its two constraints. Everything must be discrete, and everything must be made of one kind of thing, a unit of experience called a vibe. These two demands, applied without exception, do most of the work.

A single thesis organizes the whole paper. Experience is not the top of a stack that emerges from a non-experiential computational base. Experience is the bottom of the stack, the substance everything is made of. The discrete mathematics that describes how it updates, of the kind explored by Wolfram (2020), Fredkin (2003), and 't Hooft (2016), is a description of its form, layered on top, not a meaningless machine underneath. The bits are felt tones. The computation is experience computing itself.

This work is a philosophical and mathematical model rather than a finished scientific theory. It does not claim to derive or replace established physical law, and references to physics, computation, and mathematics are sources of structure and analogy unless stated otherwise. What it offers is a candidate ontology with a concrete build target, open to revision, simulation, and critique.
